\chapter*{Abstract}
\addcontentsline{toc}{chapter}{Abstract}

Two identical cells exposed to identical environmental conditions might be expected to respond identically, but the promoter controlling a gene switches on and off at random, so a cell may make a lot of protein in a burst or lie dormant instead. In HIV, these bursts can be the difference between a cell actively making virus and going latent, and because antivirals can only block replicating virus, latent cells survive treatment and can burst again once it stops. The dataset is comprised of movies of cells that have had a GFP instruction attached to their promoter region, which allows me to track frame-by-frame luminosity. Only the brightness of each cell can be observed, and it must be used to work out the rates behind it: how often the promoter turns on, how long it stays on, and how much protein it produces while on. Since the data is only given on a pixel-by-pixel basis, each cell must be segmented and tracked and then given some value to represent its luminosity at each timestep. Cellpose is used along with an implementation of the Hungarian algorithm to identify and track cells. This yielded 479 cells present in all 72 frames from hour 1 to 19 in one well. I compared the simulator against two models before proceeding, and I similarly validated the ABC on synthetic data with known rates. While it met the coverage requirement, the intervals were too wide to be useful, implying that the estimates themselves, while accurate, were uncertain. Compared with an exact-likelihood benchmark, this suggests the information was in the data and rejection ABC lost it; regression adjustment recovered it, and coverage tests confirmed the intervals stayed calibrated. On the real well, a large-number telegraph model that gives each cell its own brightness finds a burst frequency of 0.48 [0.08, 2.3] per hour, in line with the Dar lab's own estimate of 0.42, and an extrinsic spread $\eta = 0.59$ [0.53, 0.66], so most of the variation is between cells. The model reproduces the population statistics but not individual cells: real cells drift slowly where the model expects them to revert, so simply assuming each cell stays where it was makes better forecasts.
